\chapter{Artificial Neural Networks with Keras}\label{n:ann}
\studytime{2.5 hours}

\section{From Biological to Artificial Neurons}
\begin{dfn}{Artificial neural network (ANN)}
An \textbf{artificial neural network} is a machine-learning model inspired by the network of biological neurons in the brain. It consists of layers of simple connected units (\textbf{artificial neurons}); each connection has a \textbf{weight} that is learned from data. ANNs are the core of \textbf{deep learning}.
\end{dfn}
\begin{center}
\small
\begin{tabular}{@{}ll@{}}
\toprule
\textbf{Biological neuron} & \textbf{Artificial neuron}\\
\midrule
dendrites (receive signals) & inputs $x_j$\\
synapse strength & weights $w_j$\\
cell body (sums signals) & weighted sum $z=\vect w\cdot\vect x+b$\\
axon fires if above threshold & activation $a=\varphi(z)$\\
\bottomrule
\end{tabular}
\end{center}
\textbf{Why now:} huge data, powerful GPUs, better algorithms (ReLU, Adam, dropout), better initialisation.

\begin{drawit}{Artificial neuron}
\centering
\begin{tikzpicture}[font=\small]
\foreach \i/\y in {1/1.2,2/0,3/-1.2}{\node[draw,circle,fill=mlblue!10] (x\i) at (0,\y) {$x_\i$};}
\node[draw,circle,thick,fill=mlorange!15,minimum size=1.2cm] (s) at (3,0) {$\Sigma\ |\ \varphi$};
\foreach \i in {1,2,3}{\draw[-Stealth] (x\i)--(s) node[midway,above,font=\scriptsize]{$w_\i$};}
\node (b) at (3,1.6) {$b$}; \draw[-Stealth] (b)--(s);
\draw[-Stealth,thick] (s)--(5.6,0) node[right]{$a=\varphi(\vect w\cdot\vect x+b)$};
\end{tikzpicture}
\end{drawit}

\section{The Perceptron}
\begin{dfn}{Perceptron (Rosenblatt, 1957)}
The \textbf{perceptron} is the simplest ANN: a single layer of \textbf{threshold logic units (TLUs)}. Each TLU computes $z=\vect w\cdot\vect x+b$ and outputs a \textbf{step function}: $\hat y=1$ if $z>0$, else $0$. It is a linear binary classifier.
\end{dfn}
\textbf{Perceptron learning rule:} for each training instance, $\ w_j\leftarrow w_j+\eta(y-\hat y)x_j$, $\ b\leftarrow b+\eta(y-\hat y)$. Weights change only on mistakes.
\textbf{Perceptron convergence theorem:} if the data is linearly separable, the rule converges to a solution.

\begin{numer}[: perceptron learns AND ($\eta=1$, start $\vect w=(0,0)$, $b=0$)]
Epoch 1: $(1,1)$ predicted 0, $y=1$ $\Rightarrow$ $\vect w=(1,1)$, $b=1$.\\
Continuing, after epoch 5: $\vect w=(2,1)$, $b=-2$; epoch 6 makes no mistakes.\\
Check: $(1,1)\to1>0\Rightarrow1$; $(1,0)\to0\Rightarrow0$; $(0,1)\to-1\Rightarrow0$; $(0,0)\to-2\Rightarrow0$ \checkmark. Boundary $2x_1+x_2=2$.
\end{numer}

\begin{itemize}
  \item \textbf{XOR problem} (Minsky \& Papert, 1969): XOR is not linearly separable, so a single perceptron cannot learn it. Solution: \textbf{stack} perceptrons $\Rightarrow$ multilayer perceptron. E.g.\ XOR $=$ AND(OR, NAND).
\end{itemize}

\section{Multilayer Perceptron (MLP) and Backpropagation}
\begin{dfn}{Multilayer perceptron}
An \textbf{MLP} consists of an \textbf{input layer}, one or more \textbf{hidden layers} of neurons, and an \textbf{output layer}; every neuron is connected to all neurons of the next layer (\textbf{dense / fully connected}). Signals flow only forward (\textbf{feedforward network}). An ANN with many hidden layers is a \textbf{deep neural network (DNN)}.
\end{dfn}
\begin{drawit}{MLP with one hidden layer}
\centering
\begin{tikzpicture}[n/.style={circle,draw,thick,minimum size=0.6cm}]
\foreach \i in {1,2,3}{\node[n,fill=mlblue!10] (i\i) at (0,1.5-\i) {};}
\foreach \i in {1,...,4}{\node[n,fill=mlorange!15] (h\i) at (2.5,2-\i) {};}
\foreach \i in {1,2}{\node[n,fill=mlgreen!15] (o\i) at (5,1-\i) {};}
\foreach \a in {1,2,3}{\foreach \b in {1,...,4}{\draw[gray] (i\a)--(h\b);}}
\foreach \a in {1,...,4}{\foreach \b in {1,2}{\draw[gray] (h\a)--(o\b);}}
\node[font=\scriptsize] at (0,-2.6) {input}; \node[font=\scriptsize] at (2.5,-2.6) {hidden}; \node[font=\scriptsize] at (5,-2.6) {output};
\end{tikzpicture}
\end{drawit}

\begin{numer}[: forward pass of a 2--2--1 network (sigmoid)]
$\vect x=(1,0.5)$; hidden: $z_1=0.4(1)-0.6(0.5)+0.1=0.2$, $a_1=\sigma(0.2)=0.5498$; $z_2=0.7(1)+0.2(0.5)-0.2=0.6$, $a_2=0.6457$.\\
Output: $z=0.5(0.5498)-0.3(0.6457)+0.05=0.1312$, $\hat y=\sigma(0.1312)=\mathbf{0.533}$. Parameters $=(2\cdot2+2)+(2+1)=9$.
\end{numer}
\textbf{Parameters of a dense layer} $=n_{\text{in}}n_{\text{out}}+n_{\text{out}}$. MLP 784--300--100--10: $235\,500+30\,100+1\,010=\mathbf{266\,610}$.

\begin{dfn}{Backpropagation}
\textbf{Backpropagation} is the algorithm used to train MLPs. For each mini-batch it makes a \textbf{forward pass} (compute outputs and loss), then a \textbf{backward pass} that uses the \textbf{chain rule} to compute the gradient of the loss with respect to every weight, layer by layer from output to input, and finally a \textbf{gradient descent step} to update the weights.
\end{dfn}
\textbf{Steps:}
\begin{enumerate}
  \item Initialise weights randomly (not zero: symmetry breaking).
  \item Forward pass: compute every layer's output and the loss.
  \item Output error: $\delta^{(L)}=\frac{\partial J}{\partial a^{(L)}}\odot\varphi'(z^{(L)})$.
  \item Propagate back: $\delta^{(l)}=(\mathbf W^{(l+1)\top}\delta^{(l+1)})\odot\varphi'(z^{(l)})$.
  \item Gradients: $\frac{\partial J}{\partial\mathbf W^{(l)}}=\delta^{(l)}a^{(l-1)\top}$; update $\mathbf W\leftarrow\mathbf W-\eta\frac{\partial J}{\partial\mathbf W}$.
  \item Repeat for all mini-batches; one pass over the training set is an \textbf{epoch}.
\end{enumerate}
Step function has zero gradient, so MLPs use smooth activations.

\section{Activation Functions}
\begin{center}
\small
\begin{tabular}{@{}llll@{}}
\toprule
\textbf{Function} & \textbf{Formula} & \textbf{Range} & \textbf{Note}\\
\midrule
Sigmoid & $\frac1{1+e^{-z}}$ & $(0,1)$ & binary output; vanishing gradient\\
Tanh & $\frac{e^z-e^{-z}}{e^z+e^{-z}}$ & $(-1,1)$ & zero-centred\\
ReLU & $\max(0,z)$ & $[0,\infty)$ & default for hidden layers; fast; ``dying ReLU''\\
Leaky ReLU & $\max(\alpha z,z)$ & $(-\infty,\infty)$ & fixes dying ReLU\\
Softmax & $\frac{e^{z_k}}{\sum_je^{z_j}}$ & $(0,1)$, sum 1 & multi-class output\\
\bottomrule
\end{tabular}
\end{center}
\textbf{Why non-linear?} Without non-linear activations, many layers collapse into a single linear transformation and cannot learn complex patterns.

\begin{center}
\small
\begin{tabular}{@{}p{3.0cm}p{3.2cm}p{3.2cm}p{5.5cm}@{}}
\toprule
\textbf{Task} & \textbf{Output neurons} & \textbf{Output activation} & \textbf{Loss}\\
\midrule
Regression & 1 per value & none (linear) & MSE (or MAE, Huber)\\
Binary classification & 1 & sigmoid & binary cross-entropy\\
Multi-label & 1 per label & sigmoid & binary cross-entropy\\
Multi-class & 1 per class & softmax & (sparse) categorical cross-entropy\\
\bottomrule
\end{tabular}
\end{center}

\section{Implementing MLPs with Keras}
\begin{dfn}{Keras}
\textbf{Keras} is a high-level deep-learning API (part of TensorFlow as \texttt{tf.keras}) for building, training, evaluating and running neural networks easily.
\end{dfn}
\textbf{Workflow (5 steps):}
\begin{enumerate}
  \item \textbf{Load and prepare data:} split train/validation/test; scale inputs (e.g.\ pixels $/255$).
  \item \textbf{Build the model:} \texttt{Sequential} API (stack of layers: \texttt{Flatten}, \texttt{Dense(300, activation="relu")}, \dots, \texttt{Dense(10, activation="softmax")}).
  \item \textbf{Compile:} choose loss (\texttt{sparse\_categorical\_crossentropy}), optimizer (\texttt{sgd}, \texttt{adam}) and metrics (\texttt{accuracy}).
  \item \textbf{Train:} \texttt{fit(X, y, epochs, batch\_size, validation\_data)}; returns a \texttt{History} (learning curves).
  \item \textbf{Evaluate and predict:} \texttt{evaluate()} on the test set, \texttt{predict()} for new data; \texttt{save()} the model.
\end{enumerate}
\begin{itemize}
  \item \textbf{Three APIs:} \textbf{Sequential} (simple stack), \textbf{Functional} (multiple inputs/outputs, wide-and-deep networks), \textbf{Subclassing} (fully dynamic models).
  \item \textbf{Callbacks:} \texttt{ModelCheckpoint} (save best model), \texttt{EarlyStopping} (stop when validation loss stops improving), \texttt{TensorBoard} (visualise).
  \item Updates per epoch $=\lceil n/\text{batch size}\rceil$; e.g.\ $55\,000/32\to1719$.
  \item Example: Fashion-MNIST MLP reaches about 88\% test accuracy.
\end{itemize}

\section{Training Deep Networks: Key Ideas}
\begin{itemize}
  \item \textbf{Hyperparameters:} number of hidden layers, neurons per layer, learning rate (most important), batch size, optimizer, activation, epochs. Tune with validation data or random search (Keras Tuner).
  \item \textbf{Vanishing/exploding gradients:} gradients shrink or grow through many layers. Fixes: ReLU, He/Glorot initialisation, batch normalisation, gradient clipping.
  \item \textbf{Optimizers:} SGD, momentum ($\vect v\leftarrow\beta\vect v-\eta\vect g$), RMSProp, \textbf{Adam} (momentum + adaptive rates; default choice).
  \item \textbf{Regularisation (avoid overfitting):} L1/L2 weight penalty, \textbf{dropout} (randomly switch off neurons during training, e.g.\ 20--50\%), \textbf{early stopping}, more data / data augmentation.
\end{itemize}
\apa{learns complex non-linear patterns, automatic feature learning, scales with data and GPUs, state of the art for images, speech and text}{needs large data and compute, black box, many hyperparameters, can overfit}{image recognition, speech recognition, translation, recommendation, self-driving cars, medical imaging}

\begin{imp}
Most asked: perceptron and its learning rule; XOR problem; MLP architecture; backpropagation steps; activation functions table; Keras workflow (build--compile--fit--evaluate); dropout and early stopping.
\end{imp}

\begin{iq}
\item \twom What is a perceptron?
\item \twom Why can a perceptron not solve XOR?
\item \twom Define an activation function. Why is ReLU popular?
\item \twom What is an epoch?
\item \fivem Compare biological and artificial neurons.
\item \fivem Explain activation functions with formulas.
\item \fivem Explain the Sequential and Functional APIs of Keras.
\item \fivem What is dropout? How does it reduce overfitting?
\item \tenm Explain the multilayer perceptron and the backpropagation algorithm.
\item \tenm Explain the steps to build, compile, train and evaluate an MLP classifier using Keras.
\end{iq}
